\scanpage{4}
\mspage{IV.2}
if and only if $\unsure{\mathbb{L}|\mathscr{L}}$ is equivalent to $\handwritten{\mathscr{A}}$.

\uline{Proof:} Certainly if $\mathscr{L}$ is in $\handwritten{\mathscr{C}}_1(\mathrm{X},\handwritten{\mathscr{A}})$ then $\mathscr{L}$ is equivalent to $\handwritten{\mathscr{A}}$. Suppose $\mathscr{L}$ is equivalent to $\handwritten{\mathscr{A}}$. Then there are faithful $\handwritten{\mathscr{A}}$-modules $\handwritten{\mathscr{M}}_1$ and $\handwritten{\mathscr{M}}_2$, locally projective of finite type, such that
\[
\mathscr{L}\otimes_{\handwritten{\mathscr{A}}}\underline{\mathrm{Hom}}_{\handwritten{\mathscr{A}}}(\handwritten{\mathscr{M}}_1,\handwritten{\mathscr{M}}_1)\approx \underline{\mathrm{Hom}}_{\handwritten{\mathscr{A}}}(\handwritten{\mathscr{M}}_2,\handwritten{\mathscr{M}}_2).
\]

Now $\underline{\mathrm{Hom}}_{\handwritten{\mathscr{A}}}(\handwritten{\mathscr{M}}_1,\handwritten{\mathscr{M}}_1)$ is imbedded in $\mathscr{L}\otimes_{\handwritten{\mathscr{A}}}\underline{\mathrm{Hom}}_{\handwritten{\mathscr{A}}}(\handwritten{\mathscr{M}}_1,\handwritten{\mathscr{M}}_1)$, so $\handwritten{\mathscr{M}}_2$ is a module over $\underline{\mathrm{Hom}}_{\handwritten{\mathscr{A}}}(\handwritten{\mathscr{M}}_1,\handwritten{\mathscr{M}}_1)$, which we will call $\handwritten{\mathcal{M}}$.

Now we have a map of presheaves $\alpha'\colon \underline{\mathrm{Hom}}_{\handwritten{\mathcal{M}}}(\handwritten{\mathscr{M}}_1,\handwritten{\mathscr{M}}_2)\otimes_{\handwritten{\mathscr{A}}}^{\handwritten{\prime}}\handwritten{\mathscr{M}}_1\to \handwritten{\mathscr{M}}_2$, defined by
\[
\alpha'_{\mathrm U}(\{\mathrm{g}_{\mathrm V}\}_{\mathrm V\subset\mathrm U}\}\otimes\mathrm{m}_1)=\mathrm{g}_{\mathrm U}(\mathrm{m}_1)
\]
for $\{\mathrm{g}_{\mathrm V}\}_{\mathrm V\subset\mathrm U}$ in $\mathrm{Hom}_{\handwritten{\mathcal{M}}|\mathrm U}(\handwritten{\mathscr{M}}_1|\mathrm U,\handwritten{\mathscr{M}}_2|\mathrm U)$ and $\mathrm{m}_1$ in $\handwritten{\mathscr{M}}_1(\mathrm U)$. Since $\handwritten{\mathscr{M}}_1$ has finite presentation over $\handwritten{\mathcal{M}}$ (Theorem II.7), we know from [5, Theorem A.6]that the induced map of sheaves, $\alpha$, is an isomorphism. Let $\underline{\mathrm{Hom}}_{\handwritten{\mathcal{M}}}(\handwritten{\mathscr{M}}_1,\handwritten{\mathscr{M}}_2)=\handwritten{\mathscr{M}}_3$. Since $\handwritten{\mathscr{M}}_1$ is locally projective of finite type as a $\handwritten{\mathcal{M}}$-module (by Theorem II.7) and $\handwritten{\mathscr{M}}_2$ is of finite type as a $\handwritten{\mathcal{M}}$-module, $\handwritten{\mathscr{M}}_3$ is of finite type as a $\handwritten{\mathcal{M}}$-module, and therefore is of finite type as an $\handwritten{\mathscr{A}}$-module. Thus, from Theorem II. 9, $\handwritten{\mathscr{M}}_3$ is locally projective of finite type
% archive sheet number: 3
